%!TEX TS-program = latex
%!TEX encoding = UTF-8 Unicode

\documentclass[a4wide,12pt]{report} %report
\usepackage[left=1in, right=1in, top=0.75in, bottom=0.75in, includehead, headheight=17.6pt]{geometry}

%\oddsidemargin=.1in
%\evensidemargin=.2in
%\textwidth=6.3in
%\topmargin=-.5in
%\textheight=9in

\textfloatsep 0.25cm
\intextsep 0.25cm
\belowcaptionskip -0.35cm

\usepackage[hang, small, bf, margin=20pt, tableposition=top]{caption}
\setlength{\abovecaptionskip}{2pt}

\usepackage{amsmath}
\usepackage{amsfonts}
\usepackage{amssymb}
\usepackage{latexsym}
\usepackage{textcomp}
\usepackage{graphicx}

\usepackage{afterpage}
\usepackage[superscript]{cite}
\usepackage{citesort}
\usepackage{overcite}

%\usepackage{chapterbib}

\usepackage{array}
\usepackage{float}
\usepackage{fancyhdr}

%\usepackage{palatino}

%\oddsidemargin -0.35in
%\evensidemargin -.5in

%\textwidth=7in
%\textheight=9in
%\sectionskip=0.075in

%\topmargin=-0.75in
\footskip=0.35in

%\font\1 "TUE Scala" at 11pt
%\font\a "TUE Scala Bold" at 12pt
%\font\2 "TUE Scala Italic" at 14pt
%\font\7 "Lucida Bright" at 10pt
%\font\3 "Hoefler Text:mapping=tex-text;Style Options=Engraved Text;
             % Letter Case=All Capitals;color=229966" at 24pt
%\font\4 "Hoefler Text:mapping=tex-text;Letter Case=All Capitals" at 9 pt

%\font\5="Apple Chancery:Letter Case=Small Caps;
%Design Complexity=Simple Design Level" at 10pt
%\font\6="Hoefler Text:%Letter Case=Small Caps;
%Design Complexity=Simple Design Level" at 10pt
%\font\6="TUE Scala:Letter Case=Small Caps;Design
%Complexity=Simple Design Level" at 12pt
\usepackage{setspace} 
\doublespacing

\title{Polymeric Building Blocks Prepared by RAFT Polymerization for Supramolecular Assembly Based on Cucurbit[8]uril}
%\date{}
\setcounter{page}{1}
\setcounter{tocdepth}{1}
%\setcounter{secnumdepth}{-2}


\pagestyle{fancy}
%\lhead{\it Research Proposal SCR}
%\rhead{Eric Andrew Appel}
%Melville Laboratory for Polymer Synthesis\\
%University of Cambridge}

%\fancyfoot[C]{page \thepage}
\cfoot{\thepage}



%%%%%%%%%%%%%%%%%%%%%%%%%%%%%%%%%%%%%%%%%%%%%%%%%%%%%%%%%%%%%%%%%%%%%%%%%%%



\begin{document}
\begin{titlepage}


\begin{figure}[h]
%\begin{center}
\begin{flushright}
\includegraphics[width=0.45\textwidth]{Images/CUnibig.eps}
\end{flushright}
\end{figure}
\begin{flushright}
\text{Melville Laboratory for Polymer Synthesis} \\
\text{The University Chemical Laboratory} \\
\text{Lensfield Road, Cambridge, CB2 1EW} \\
\text{United Kingdom} \\
\end{flushright}
%\end{center}



\vspace*{0.5cm}
\begin{center}
\Huge {Polymeric Building Blocks Prepared by RAFT Polymerization for Supramolecular Assembly Based on Cucurbit[8]uril}
\end{center}
\vspace{0.75cm}
\begin{center}
\Large {Eric Andrew Appel} \\[3pt]
%\textit{Jesus College}\\[1cm]
\end{center}


\begin{figure}[h]
\begin{center}
\includegraphics[width=0.15\textwidth]{Images/jesus.eps}
\end{center}
\end{figure}

\vspace{-1.2cm}
\begin{center}
\Large {\it Jesus College}\\
\end{center}

\vspace{0.25cm}
\begin{center}
\Large {Dr Oren A. Scherman}\\
\today
\end{center}


\end{titlepage}



%%%%%%%%%%%%%%%%%%%%%%%%%%%%%%%%%%%%%%%%%%%%%%%%%%%%%%%%%%%%%%%%%%%%%%%%%%%



\doublespacing


\begin{center}
\Large{Declaration}
\end{center}
\noindent
This report is submitted in partial fulfillment of the requirements for the Certificate of Post-Graduate Study in Chemistry. Except where indicated to the contrary, either directly or by reference, the work described in this dissertation is solely the work of the author. 
\\
\\ 
\\
\begin{flushright}
Eric Andrew Appel\\
University of Cambridge\\
\today
\end{flushright}


\newpage
\begin{center}
\Large{Acknowledgements}
\end{center}
\noindent
I would first like to thank my supervisor Dr. Oren A. Scherman for a very interesting and stimulating research topic, his guidance and attitude and especially for searching diligently for funding for this project to bring me out from the US. I would also like to thank my fellow coworkers in the Scherman group for making this first year here very enjoyable, especially Urs, Frank, Jameel, Adam, Nicolas and Monika for their help in initially finding my feet here. I would also like to acknowledge everyone else in the Melville Laboratory for Polymer Synthesis for a very friendly, helpful and insightful working environment. I would like to thank Trevor L. Hughes and Matthew J. Miller from Schlumberger for funding and support. Finally, I would like to thank my wife for feeding me well and keeping me sane.


\newpage
\begin{center}
\Large{Abbreviations}
\end{center}
\noindent
AA - acrylic acid\\
ACPA - azobiscyanopentanoic acid\\
AIBN - azobisisobutyronitrile\\
AM - acrylamide\\
AN - acrylonitrile\\
ATRP - atom transfer radical polymerization\\
CB[n] - cucurbit[n]uril\\
{\it x}-CD - cyclodextrin with {\it x} glucopyranoside units ({\it x} = $\alpha$ (6), $\beta$ (7) and $\gamma$ (8))\\
CDCl$_3$ - deuterochloroform\\
DCC - dicyclohexylcarbodiimide\\
DCM - dichloromethane\\
DDMAT - S-dodecyl-S'-($\alpha$,$\alpha$'-dimethyl-$\alpha$''-acetic acid) trithiocarbonate\\
DMA - N,N-dimethylacrylamide\\
DMF - dimethylformamide\\
D$_2$O - deuterated water\\
DMSO - dimethylsulfoxide\\
EA - ethyl acetate\\
H-bond - hydrogen bond\\
HOBt - N-hydroxybenzotriazole\\
HTMPB - hexadecyltrimethylphosphonium bromide\\
%ITC - isothermal titration calorimetry\\
K$_a$ - association constant / M$^{-1}$\\
K$_d$ - dissociation constant / M$^{1}$\\
MA - methyl acrylate\\
MADIX - macromolecular design {\it via} the interchange of xanthates\\
MBipy - 1-methyl-4,4$'$-bipyridinium iodide\\
MeOD - deuterated methanol\\
MHGC - macrocyclic host-guest complexation\\
MLC - metal-ligand complexation\\
MMA - methyl methacrylate\\
M2V - viologen (doubly alkylated 4,4-bypyridine moiety)\\
Napy - 2,7-diamido-1,8-naphthyridine\\
NIPAM - N-isopropylacrylamide\\
Np - 2-naphthol\\
NVC - N-vinylcarbozole\\
NVP - N-vinylpyrrolidone\\
PEB - poly(ethylene-co-butene)\\
PEG - poly(ethylene glycol)\\
RAFT - reversible addition-fragmentation chain-transfer\\
ROMP - ring-opening metathesis polymerization\\
ROP - ring-opening polymerization\\
S - styrene\\
SS - styrene sulfonate\\
tBA - poly({\it tert}-butylacrylate)\\
TEG - triethyleneglycol\\
THF - tetrahydrofuran\\
UPy - 2-ureido-4[1{\it H}]-pyrimidinone\\
VAc - vinyl acetate\\
{\it x}-VP - {\it x}-vinylpyridine\\
%VL - $\delta$-valerolactone\\
{\it y}@{\it x} - guest {\it y} encapsulated within host {\it x}


\tableofcontents

\fancyhead[RE, LO]{Eric Andrew Appel\\ \it Contents}
\fancyhead[RO, LE]{Melville Laboratory for Polymer Synthesis\\ University of Cambridge}



%%%%%%%%%%%%%%%%%%%%%%%%%%%%%%%%%%%%%%%%%%%%%%%%%%%%%%%%%%%%%%%%%%%%%%%%%%%




\chapter{Introduction}

\fancyhead[RE, LO]{Eric Andrew Appel\\ \it Chapter 1 - Introduction}
\fancyhead[RO, LE]{Melville Laboratory for Polymer Synthesis\\ University of Cambridge}


Supramolecular chemistry \cite{Lehn:1995p2342, Cram:1974p2474, PEDERSEN:1988p2572} is an extremely fast growing field, which focuses greatly on the concept of self assembly on the molecular level. 



%%%%%%%%%%%%%%%%%%%%%%%%%%%%%%%%%%
\\
\\
\section{Dynamic, Supramolecular Polymers}
\noindent
Despite extreme initial skepticism facing macromolecular science in the early 20th century, the development of polymers over the past century has had a dramatic and indispensable impact on human life and earned a Nobel Prize for Hermann Staudinger in 1953.\cite{Staudinger:1953p3048}


%%%%%%%%%%%%%%%%%%%%%%%%%%%%%%%%%%
\\
\subsection{Characteristics}
\noindent
In order to look at the general characteristics of SPs, certain limitations are necessarily required for the definition of `supramolecular polymer'. 



%%%%%%%%%%%%%%%%%%%%%%%%%%%%%%%%%%


\subsection{Non-covalent interactions and macro-SPs}
\noindent
As the characteristics of SPs are dependent almost exclusively on the non-covalent interactions existing between monomers, it is important to look at the self-assembly toolbox to identify advantages and limitations existing in each of the three major categories in the formation of SPs: (a) hydrogen bonding (H-bonding), (b) metal-ligand coordination (MLC) and (c) macrocyclic host-guest complexation (MHGC).

%%%%%%%%%%%%%%%%%%%%%%%%%%%%%%%%%%
\\
\\
\section{Cucurbit[8]uril: A Host of Opportunities}
\noindent
Although it was first synthesized more than a century ago,\cite{Behrend:1905p1004} the structure of the synthetic host molecule cucurbituril (CB) was not elucidated until the 1980s.\cite{Freeman:1981p990}

\vspace{0.25 in}
\begin{figure}[h]
\begin{center}
 \includegraphics[width=0.9\textwidth]{Images/generic_CB.eps}
\vspace{0.1 in}
  \caption{Schematic representation of the cucurbit[n]uril (CB[n]) family demonstrating the dimensions of the supramolecular hosts and cartoon depiction as a barrel with symmetric portals. A 3D model of cucurbit[8]uril is shown at right. Image courtesy of Urs Rauwald.}\label{generic_CB}
\end{center}
\vspace{0.25 in}
\end{figure} 


%%%%%%%%%%%%%%%%%%%%%%


\\
\\
\section{RAFT Polymerization}

\noindent
`Living' polymerization techniques have made enormous advances since the inception of anionic polymerization in 1956 by Szwarc {\it et.~al.}~\cite{Szwarc:1956p46, Szwarc:1956p55} 

%%%%%%%%%%%%%%%%%%%%%%%%%%%%%%%%%%%%%%%%%%%%%%%%%%%%%%%%%%%%%%%%%%%%%



\chapter{Results and Discussion}

\fancyhead[RE, LO]{Eric Andrew Appel\\ \it Chapter 2 - Results and Discussion}
\fancyhead[RO, LE]{Melville Laboratory for Polymer Synthesis\\ University of Cambridge}



\section{The Search for a Simple Path to Functional RAFT CTAs}

\noindent
Initial efforts focused on the synthesis of trithiocarbonate (TTC) CTAs useful for RAFT polymerization that carry Np or M2V functionality.  



\section{Functional Water Soluble Polymers}

With several RAFT CTAs in hand it was now possible to prepare functional polymers.



\section{Supramolecular Polymer Assemblies in Water}


As several functional polymers of various types were now in hand, it was time to determine the capacity of the polymers to complex with M2V and CB[8].


%%%%%%%%%%%%%%%%%%%%%%%%%%%%%%%%%%%%%%%%%%%%%%%%%%%%%%%%%%%%%%%%%%%%%%%%



\chapter{Conclusions}

\fancyhead[RE, LO]{Eric Andrew Appel\\ \it Chapter 3 - Conclusions}
\fancyhead[RO, LE]{Melville Laboratory for Polymer Synthesis\\ University of Cambridge}



Mono- and bis-Np functional RAFT CTAs have have been synthesized and have demonstrated to successfully confer control over the polymerization of a variety of hydrophobic, hydrophilic and stimuli-responsive monomers with excellent control over targeted molecular weight and PDI. 



%%%%%%%%%%%%%%%%%%%%%%%%%%%%%%%%%%%%%%%%%%%%%%%%%%%%%%%%%%%%%%%%%%%%%%%%



\chapter{Outlook and Future Work}

\fancyhead[RE, LO]{Eric Andrew Appel\\ \it Chapter 4 - Outlook}
\fancyhead[RO, LE]{Melville Laboratory for Polymer Synthesis\\ University of Cambridge}

The development of functional RAFT CTAs that provide facile preparation of either semi-telechelic and telechelic polymers that can be dynamically interlinked with CB[8] offers many opportunities for further investigation into hierarchical self-assembly.



%%%%%%%%%%%%%%%%%%%%%%%%%%%%%%%%%%%%%%%%%%%%%%%%%%%%%%%%%%%%%%%%%%%%%%%%



\chapter{Experimental}



\fancyhead[RE, LO]{Eric Andrew Appel\\ \it Chapter 5 - Experimental}
\fancyhead[RO, LE]{Melville Laboratory for Polymer Synthesis\\ University of Cambridge}


\section{Instrumentation and Materials}
\noindent  
$^{1}$H NMR (400~MHz) spectra was recorded using a Bruker Avance QNP~400. Chemical shifts are recorded in ppm ($\delta$) in CDCl$_{3}$ with the internal reference set to $\delta$~7.26~ppm or MeOD with the internal reference set to $\delta$~3.31~ppm.
%
 $^{13}$C NMR (125~MHz) spectra was recorded using a Bruker Avance Cryobrobe ATM TCI DRX 500 or a Bruker Avance 500 BB-ATM. Chemical shifts are recorded in ppm ($\delta$) in CDCl$_{3}$ and MeOD with the internal reference set to $\delta$~77.16~ppm and $\delta$~49.00~ppm, respectively.
%
ATR FT-IR spectroscopy was performed using a Perkin-Elmer Spectrum 100 series FT-IR spectrometer equipped with a universal ATR sampling accessory. 
%
High-resolution mass spectra was recorded on a Bruker BioASpex II 4.7e FT-ICR mass spectrometer liquid chromatography-mass spectrometry Waters ZQ. 
%
UV-VIS studies were performed on a Varian Cary 4000 UV-Vis spectrophotometer. 
%
Gel permeation chromatography (GPC) was carried out in either tetrahydrofuran (THF) or in dimethylformamide (DMF). 
%
THF GPC was performed on two Jordi DVB columns, with pore sizes ranging from 100--10$^{5}$\AA, connected in series with a SPD-M20A prominence diode array detector (Shimadzu), an Optilab DSP refractive index detector and a Dawn DSP multi-angle light scattering detector (both Wyatt technologies) calibrated with relation to polystyrene standards. 
%
Samples were filtered over 0.45~$\mu$m PTFE filters before injection using a 1.0~mL~/~min flow rate.
%
DMF GPC was performed on two Jordi 5$\mu$m DVB columns connected in series with a SPD-M20A prominence diode array detector  and refractive index detector (both Shimadzu) calibrated in relation to poly(methyl methacrylate) standards. 
%
Samples were filtered over 0.45~$\mu$m nylon filters before injection using a 0.75~mL~/~min flow rate. 
%
Solution viscosities were measured using Schott-Ger\"{a}te Ubbelohde micro-viscometers with a suspended level bulb using a PVS1 measuring device, the micro-viscometers were thermostated in a PV15 water bath at 25.00~($\pm$ 0.01)~$^{\textrm o}$C using a DLK10 thermostat unit (all manufactured by Lauda). 
%
Intrinsic viscosities were corrected using the appropriate Hagenbach correction factors. 
%
%Differential scanning calorimetry (DSC) was performed using a Perkin Elmer Pyris 1 calorimeter and CryoFill liquid nitrogen cooling system.  
\\
\\
DCM, chloroform, DMF and THF were dispensed directly from alumina based solvent columns (Puresolv) under an inert atmosphere. 
%
TEA was purchased from Sigma-Aldrich and dried over potassium hydroxide. 
%
Acrylic Acid (AA), t-butyl acrylate (t-BA), and N-N-Dimethylacrylamide (DMA) were purchased from Aldrich and purified on a basic alumina column to remove the inhibitor and stored over {4\AA} molecular sieves.
%
N-isopropyl acrylamide (NIPAM) (Sigma-Aldrich) was recrystallized twice from hexane and dried under vacuum. 
%
Acrylamide (AM) (Sigma-Aldrich) was used as received.
%
2,2′-Azobis(2-methylpropionitrile) (Sigma-Aldrich) and 4,4'-Azobis(4-cyanovaleric acid) (Sigma-Aldrich) were recrystallized twice from methanol. 
%
2-(4-(bromomethyl)benzyloxy)naphthalene {\bf 24} was prepared according to literature procedure.\cite{Kim:2004p1558}
%
2-(dodecanethiocarbonothioylthio)-2-methylpropanoic acid (DDMAT) {\bf 26} was prepared according to literature procedure.\cite{Skey:2008p4456}
%
1-(2-aminoethyl)-1'-methyl-4,4'-bipyridine-1,1'-diium hexafluorophosphate(V) {\bf 31} was provided by Frank Biederman.
%
All other materials were purchased from Sigma-Aldrich and used as received.





\section{Synthetic Methods}


\vspace{.3 in}
\subsection{Synthesis of Np functional CTA {\bf 25}}
%
2-hydroxyethyl 4-((naphthalen-2-yloxy)methyl)benzyl carbonotrithioate {\bf 25} was prepared by adding 2-mercaptoethanol (0.1~g, 1.3~mmol) to a suspension of K$_3$PO$_4$ (0.54~g, 2.5~mmol) in acetone (20~mL) and stirring for 30 min. 
%
CS$_2$ (0.30~g, 4.0~mmol) was added {\it via} syringe and the solution turned bright yellow after 5~min.
%
A solution of {\bf 24} (0.5~g, 1.5~mmol) in acetone (10~mL) was then added {\it via} syringe and the reaction mixture was allowed to stir for 5~h at room temperature. 
%
The reaction mixture was filtered and the cake washed with acetone (2 x 20~mL). 
%
The solvent was removed from the filtrate and the crude oil was purified by recrystallization twice from ethyl acetate/petroleum ether (40:60) to afford the title compound {\bf 25} as a yellow crystalline solid (0.42~g, 41~\%). 
%
$^1$H-NMR Spectroscopy (CDCl$_3$, 400 MHz) $\delta$ (ppm) = 7.79-7.69 (3H, m, Np-{\it H}), 7.46-7.30 (6H, m, Np-{\it H} and Ar-{\it H}), 7.24-7.19 (2H, m, Ar-{\it H}), 5.17 (2H, s, Np-O-C{\it H}$_2$-Ar), 4.64 (2H, s, Ar-C{\it H}$_2$-S), 3.90 (2H, t, HO-C{\it H}$_2$-CH$_2$-S), 3.63 (2H, t, HO-CH$_2$-C{\it H}$_2$-S).
%
$^{13}$C-NMR Spectroscopy (CDCl$_3$, 125 MHz) $\delta$ (ppm) = 223.4, 156.7, 136.7, 134.8, 134.5, 129.6, 129.6, 129.2, 128.0, 127.7, 126.9, 126.5, 123.8, 119.0, 107.2, 69.7, 60.7, 41.4, 39.3.
%
%FT-IR (ATR) $\nu$ = ... cm$^{-1}$.
%
Elemental: Found C, 63.20; H, 5.02; C$_{21}$H$_{20}$O$_2$S$_3$ required C, 62.96; H, 5.04.
%



%%%%%%%%%%%%%%%%%%%%%%%%%%%%%%%%%%%%%%%%%%%%%%%%%%%%%%%%%%%%%%%%%%%%%%%%








\singlespacing
%\restepcounter{chapter}
\addcontentsline{toc}{chapter}{\bibname}

\bibliographystyle{achemso}

\fancyhead[RE, LO]{Eric Andrew Appel\\ \it Bibliography}
\fancyhead[RO, LE]{Melville Laboratory for Polymer Synthesis\\ University of Cambridge}

\bibliography{cpgs}




\end{document}
